%% --------------------------------------------------------------------
%% Template  UTEC Tesis - Escuela de Computación
%% --------------------------------------------------------------------
% Este template ha sido modificado y adaptado para la escuela de Computación en base al template trabajo por Eduardo Castro y Roosevelt Ubaldo.
% Última actualización: Agosto, 2022

\documentclass[a4paper,12pt,oneside]{tesisutec}

\selectlanguage{spanish}

%% Paquetes
\usepackage[utf8]{inputenc}
\usepackage[square, numbers, comma, sort&compress]{natbib}
% Libreria de idioma
\usepackage[spanish]{babel}
% Libreria para posicionamiento
\usepackage{float}
% Librerias para insertar codigos
\usepackage[spanish,onelanguage,ruled,vlined]{algorithm2e}
\usepackage{verbatim} 
% Librería para hipervínculos
\usepackage{hyperref}
 % Librería necesaria para arreglar el orden de referencias en overleaf.com
\usepackage{notoccite}
\usepackage{comment}
% Diagramas vectoriales (TikZ)
\usepackage{tikz}
\usetikzlibrary{arrows.meta, positioning, calc, decorations.pathreplacing}
\usepackage{pgfplots}
\pgfplotsset{compat=1.18}

% Incluir acá los paquetes adicionales que deseas. 
% Ubicación de los imágenes.
\graphicspath{images/}

\makeatletter
% Reinsert missing \algbackskip
\def\algbackskip{\hskip-\ALG@thistlm}
\makeatother

\hypersetup{urlcolor=blue, colorlinks=true}

\begin{document}

\frontmatter

\department{CIENCIA DE LA COMPUTACIÓN}
\degree{Ciencia de la Computación}
\major{Ciencia de la Computación}

% TODO(humano): confirmar título final, ORCID y asesor antes de la entrega.
\title{CAVA: SEGMENTACIÓN TEMÁTICA DE VIDEO LARGO MEDIANTE AGREGACIÓN MULTIMODAL POR CONSENSO, SIN FUSIÓN DE REPRESENTACIONES}

\author{Edgar Moises Chambilla Mamani \orcid{0000-0000-0000-0000} \\ Sebastian Andres Quispe Barriga \orcid{0000-0000-0000-0000}} % Es obligatorio agregar ORCID del alumno
\supervisor{Edward Jorge Yuri Cayllahua Cahuina \orcid{0000-0000-0000-0000}} % Es obligatorio agregar ORCID del asesor

\date{2026}

\maketitle

\setstretch{1.5}


%% ============================================================================
%  Dedicatoria
%% ============================================================================

\dedicatory{
xxxxxxxxxxxxxxxxxxxxxxxxxxxxxxxxxxxxxxxxxxxxxxxxxxx
xxxxxxxxxxxxxxxxxxxxxxxxxxxxxxxxxxxxxxxxxxxxxxxxxxx
xxxxxxxxxxxxxxxxxxxxxxxxxxxxxxxxxxxxxxxxxxxxxxxxxxx
xxxxxxxxxxxxxxxxxxxxxxxxxxxxxxxxxxxxxxxxxxxxxxxxxxx
xxxxxxxxxxxxxxxxxxxxxxxxxxxxxxxxxxxxxxxxxxxxxxxxxxx
XXxxxxxxxxxxxxxxxxxxxxxxxxxxxxxxxxxxxxxxxxxxxxxxxx
}

%% ============================================================================
%  Pagina de Agradecimientos
%% ============================================================================

\acknowledgements{

xxxxxxxxxxxxxxxxxxxxxxxxxxxxxxxxxxxxxxxxxxxxxxxxxxx
xxxxxxxxxxxxxxxxxxxxxxxxxxxxxxxxxxxxxxxxxxxxxxxxxxx
xxxxxxxxxxxxxxxxxxxxxxxxxxxxxxxxxxxxxxxxxxxxxxxxxxx
xxxxxxxxxxxxxxxxxxxxxxxxxxxxxxxxxxxxxxxxxxxxxxxxxxx
xxxxxxxxxxxxxxxxxxxxxxxxxxxxxxxxxxxxxxxxxxxxxxxxxxx
xxxxxxxxxxxxxxxxxxxxxxxxxxxxxxxxxxxxxxxxxxxxxxxxxxx

}



\tableofcontents

\newpage

\listoftables

\newpage

\listoffigures

\addtocontents{toc}{\vspace{1.5em}}

%% ============================================================================
\mainmatter
\pagestyle{fancy}

\customchapter{RESUMEN}

Para analizar video largo de forma automática hay que dividirlo en unidades con sentido
propio, y ninguna modalidad debería ocultar a las demás al hacerlo. Esta tesis parte de un
diagnóstico de trabajo previo: en A2Summ-X, una extensión del modelo de resumen A2Summ
\cite{he2023a2summ}, la fusión de modalidades en una representación única hace que el texto
domine la ponderación entre modalidades (pesos de 0.75 a 0.89) y anule la contribución visual
antes de cualquier decisión. Se propone CAVA (\emph{Consumer-Adaptive Video Architecture}),
un sistema de segmentación temática de video largo que no fusiona representaciones. Cada
señal (transcripción, imagen, voz y las descripciones de un modelo visión--lenguaje) emite
sus propias fronteras y etiquetas de tópico. El tópico surge de un consenso de votos con
quórum, las señales sin evidencia en un tramo se abstienen, y cada grupo registra qué
señales lo sostienen.

El sistema produce un mapa de tópicos en dos niveles de granularidad que enlaza segmentos no
adyacentes cuando un tema se retoma. A partir del mapa se generan cinco formatos de versión:
capítulos, recapitulación de 60 segundos, destacados del 10\,\%, hilos de tópico e hilos por
hablante. La validación usa compuertas pre-registradas sobre datos no vistos. En las 47
películas del split de test de MovieNet \cite{huang2020movienet}, el mapa con abstención
supera a la fusión por concatenación ($z = 2.46$) y al gestor contiguo ($z = 2.48$) en MCC
por pares. La ganancia depende de la abstención: sin ella, el consenso no se distingue de la
fusión ($z = 0.80$), y la señal visual sola obtiene un MCC más alto que cualquier
combinación. Las versiones cumplen su presupuesto en todos los tipos de contenido evaluados
(película, reunión, video web y de usuario), y la identidad de hablante alcanza un ARI de 0.73 en
AMI \cite{carletta2007ami}. Tres compuertas no se cumplen: la calidad de los destacados
frente a bases simples (EXP-031/037/038), los enlaces no contiguos juzgados como la misma
trama (G-enlace-2) y la detección del narrador en documentales (EXP-039). De los 26 módulos
del sistema congelado, 10 están validados. La evaluación con personas (EXP-040) está
diseñada, con el conjunto de 9 videos propios ya procesado, y queda pendiente de
pre-registro y ejecución.

\vspace{0.5em}
\noindent \textbf{Palabras clave:}\\
\noindent Segmentación temática de video; agregación multimodal; consenso de votos;
modelos visión--lenguaje; versiones adaptadas al consumidor.

\customchapter{ABSTRACT}

\begin{center}
\large \vspace{-1.5cm} \textbf{CAVA: TOPIC SEGMENTATION OF LONG VIDEO BY CONSENSUS-BASED MULTIMODAL AGGREGATION WITHOUT REPRESENTATION FUSION}
\end{center}

Automatic analysis of long video requires splitting it into meaningful units, and no
modality should hide the others while doing so. This thesis builds on a diagnosis from
previous work: in A2Summ-X, an extension of the A2Summ summarization model
\cite{he2023a2summ}, fusing modalities into a single representation lets text dominate the
weighting between modalities (weights of 0.75 to 0.89) and cancels the visual contribution
before any decision is made. We propose CAVA (Consumer-Adaptive Video Architecture), a topic
segmentation system for long video that does not fuse representations. Each signal
(transcript, image, voice and the descriptions of a vision--language model) emits its own
boundaries and topic labels. Topics emerge from a quorum vote, signals without evidence in a
stretch abstain, and every group records which signals support it.

The system produces a topic map at two granularity levels that links non-adjacent segments
when a topic returns. Five version formats are derived from the map: chapters, a 60-second
recap, 10\,\% highlights, topic threads and speaker threads. Validation uses pre-registered
gates on unseen data. On the 47 films of the MovieNet test split \cite{huang2020movienet},
the map with abstention beats concatenation fusion ($z = 2.46$) and the contiguous manager
($z = 2.48$) in pairwise MCC. The gain depends on abstention: without it, consensus is
indistinguishable from fusion ($z = 0.80$), and the visual signal alone scores a higher MCC
than any combination. Versions meet their budget in every content type evaluated (film,
meeting, web and user video), and speaker identity reaches an ARI of 0.73 on AMI
\cite{carletta2007ami}. Three gates are not met: highlight quality against simple baselines
(EXP-031/037/038), non-adjacent links judged as the same plot (G-enlace-2), and narrator
detection in documentaries (EXP-039). Of the 26 modules of the frozen system, 10 are
validated. The human evaluation (EXP-040) is designed, with its set of 9 videos already
processed, and is pending pre-registration and execution.

\vspace{0.5em}
\noindent \textbf{Keywords:}\\
\noindent Video topic segmentation; multimodal aggregation; consensus voting;
vision--language models; consumer-adaptive versions.
 
\customchapter{INTRODUCCIÓN}

El contenido audiovisual crece más rápido de lo que una persona puede ver, y para indexarlo,
previsualizarlo o navegarlo hay que saber dónde termina cada unidad de sentido y qué
segmentos, aunque estén separados en el tiempo, tratan el mismo tema. La \emph{segmentación
temática} agrupa el contenido en tópicos coherentes. A diferencia de la segmentación en
escenas, admite que un tema se abandone y se retome más adelante, y permite que el mismo
video se consuma de formas distintas según el interés y el tiempo de quien lo ve.

Esta tesis continúa dos trabajos previos del mismo grupo. El primero, A2Summ-X, extendió el
modelo de resumen multimodal A2Summ \cite{he2023a2summ} y se evaluó con 5 splits y 5 semillas
por configuración y pruebas t pareadas sobre SumMe y TVSum \cite{gygli2014summe,
song2015tvsum}. La rama de texto aportó una mejora pequeña ($\Delta$F1 $= +0.033$,
$p = 0.014$ en TVSum), los backbones visuales modernos no mejoraron el resultado y en
conjuntos de pocos videos la varianza entre corridas dominó sobre los efectos. El segundo, SegLoop,
usó un modelo visión--lenguaje (VLM) congelado \cite{bai2025qwen25vl} como juez para generar
pseudoetiquetas con las que se entrena un segmentador de escenas ligero \cite{mun2022bassl,
huang2020movienet}. De SegLoop queda una infraestructura verificada de extracción de señales
(reconocimiento de voz, descriptores visuales y detección de cortes) que CAVA reutiliza.

El análisis que cerró A2Summ-X encontró el problema que esta tesis aborda: al fusionar las
señales en un embedding o un score único, unas reescriben a otras antes de la decisión. En
A2Summ-X, la alineación inyecta la transcripción en la consulta visual, la compuerta softmax
reparte un peso total fijo entre modalidades y la cabeza de decisión solo ve la mezcla.
Cuatro mediciones independientes coincidieron: pesos de la compuerta de texto entre 0.75 y
0.89, colapso de la señal visual en BLiSS, correlación de importancia $\rho \approx 0$ fuera
de dominio y secuencias sin cortes detectables. De ahí sale la regla de diseño de CAVA:
mantener separada la identidad de cada señal hasta el momento de decidir.

En CAVA, cada señal (transcripción, imagen, voz y las descripciones que un VLM produce por
bloques del video) calcula sus propias fronteras con su propio percentil y asigna sus propias
etiquetas de tópico. El tópico final surge de un consenso de votos con quórum. Una señal sin
evidencia en un tramo se abstiene, y cada grupo resultante registra qué señales lo sostienen
(su \emph{procedencia}). Sobre este gestor se construye el mapa de tópicos del video, con
enlaces entre segmentos no adyacentes, dos niveles de granularidad y contexto por tramo. Del
mapa se derivan cinco formatos de versión: capítulos, recapitulación de 60 segundos,
destacados del 10\,\%, hilos de tópico e hilos por hablante. El documento reporta el estado
del sistema al congelarlo: el gestor está validado en el split de test de MovieNet, las
compuertas de coherencia y de versiones están medidas, tres compuertas no se cumplen y la
evaluación con personas (EXP-040) está diseñada.

\introsection{Formulación del problema}

¿Es posible segmentar temáticamente video largo de contenido heterogéneo (película,
documental, reunión, video web) y producir versiones adaptadas al consumidor, sin anotación
humana de tópicos y sin fusionar representaciones entre señales, de modo que:
(i) el mapa de tópicos por consenso supere a la fusión de señales y a la segmentación
contigua en una métrica pre-registrada sobre datos no vistos;
(ii) los tópicos sean coherentes dentro de cada señal (margen intra$-$inter positivo);
(iii) cada versión respete su presupuesto de duración; y
(iv) un juez independiente reconozca los enlaces no contiguos como retornos del mismo tema?

\introsection{Objetivos de investigación}

\paragraph{Objetivo general}
Construir y validar experimentalmente CAVA, un sistema que segmenta temáticamente video largo
mediante agregación multimodal por consenso, sin fusión de representaciones, con abstención y
procedencia por señal, y que produce varias versiones del contenido adaptadas al consumidor,
validadas con compuertas pre-registradas sobre datos no vistos.

\paragraph{Objetivos específicos}
\begin{enumerate}
  \item Implementar las señales del sistema sobre la infraestructura de SegLoop:
        transcripción \cite{radford2023whisper} proyectada a embeddings, señal visual
        SigLIP \cite{zhai2023siglip} por unidad, voz (log-mel y tono) y descripciones de un
        VLM congelado \cite{bai2025qwen25vl}, con perfiles por tipo de contenido (película,
        documental, reunión) y ajuste de fronteras a cortes visuales.
  \item Implementar el gestor multimodal sin fusión: fronteras y etiquetas propias por
        señal, consenso de cortes por votos ponderados con tolerancia, enlace de tópicos por
        quórum (componentes conexas o enlace promedio), abstención de señales sin evidencia y
        procedencia por señal en cada grupo.
  \item Producir el mapa de tópicos con retornos, dos niveles de granularidad, etiquetas
        generadas por un modelo de lenguaje, nombres de personajes propagados por los
        enlaces y contexto por tramo, junto con la identidad de hablante (embeddings
        WavLM-sv \cite{chen2022wavlm} agrupados con Ward y BIC).
  \item Generar los cinco formatos de versión con verificación automática de presupuestos
        (compuerta G-versión).
  \item Validar el sistema con compuertas pre-registradas: el gestor frente a la fusión y al
        gestor contiguo en el split de test de MovieNet (47 películas); la coherencia
        intra$-$inter por señal; la precisión de los enlaces no contiguos con un juez ciego;
        la identidad de hablante en AMI \cite{carletta2007ami} y VoxConverse
        \cite{chung2020voxconverse}; y la calidad de los clips frente a referencias
        automáticas y bases simples.
\end{enumerate}

\introsection{Justificación}

La segmentación temática con enlaces no contiguos es necesaria para aplicaciones que la
segmentación en escenas no cubre: navegar un video por tema, generar recapitulaciones con un
tiempo fijo y reunir en un hilo todas las apariciones de un tema o de un hablante. Los
sistemas actuales dependen de VLM costosos en inferencia \cite{berman2025scenevlm} o
fusionan las modalidades antes de decidir, y en ese segundo caso el trabajo previo del grupo
mostró que una señal termina dominando a las demás. Si el consenso tardío produce mapas de
tópicos iguales o mejores que la fusión, se gana además la posibilidad de saber qué señales
sostienen cada grupo, algo que una representación fusionada no permite recuperar.

En lo metodológico, el trabajo aplica a un problema nuevo el esquema de compuertas
pre-registradas del grupo: separación de desarrollo, confirmación y test, y selección de
hiperparámetros sin consultar test. Como producto queda un pipeline reproducible (115
corridas con manifiesto que registra que no se usó anotación, 25 pruebas automatizadas y un
registro de decisiones de arquitectura) y un inventario del estado de validación de cada
módulo por tipo de contenido.

\introsection{Alcance y limitaciones}

La investigación cubre el ciclo completo, desde la extracción de señales hasta las versiones
por consumidor, con estos límites:
(i) los dominios evaluados son película (MovieNet), documental (BBC Planet Earth), reunión
(AMI) y video web y de usuario (VidChapters \cite{ventura2024vidchapters}, TVSum y SumMe),
sin anotación nueva de tópicos salvo la del conjunto de la evaluación con personas;
(ii) la calidad percibida de las versiones se mide con presupuestos y referencias
automáticas, y la preferencia humana (EXP-040) queda diseñada pero no ejecutada;
(iii) las identidades por apariencia (P\#) no se publican como identidad, y solo la identidad
por voz está validada, y únicamente en habla conversacional;
(iv) el texto en pantalla (OCR) y el espectro de audio están definidos como contratos de
señal pero no se extraen; y
(v) el cómputo se limita al clúster institucional (GPU RTX A6000 de 48\,GB y A100),
con límites de tiempo por trabajo.

\chapter{REVISIÓN CRÍTICA DE LA LITERATURA}

Este capítulo revisa las fuentes que definen el estado del arte relevante para CAVA. Empieza
por la línea de resumen de video multimodal de la que proviene este trabajo (A2Summ y
A2Summ-X), donde se originó el diagnóstico contra la fusión. Sigue con la segmentación de
escenas y capítulos (BaSSL, Scene-VLM, VidChapters), la segmentación temática de texto, la
discusión entre fusión temprana y tardía y los ensambles de agrupamientos, que son el
antecedente formal más cercano del consenso de votos. Cierra con el sistema previo SegLoop y
con el uso de VLM como jueces. Cada sección analiza los resultados reportados por sus autores
y la última resume el aporte de la tesis.

\section{Resumen de video multimodal: A2Summ}

He et al.\ \cite{he2023a2summ} proponen A2Summ, un transformer que alinea las secuencias de
video y texto y las fusiona con dos pérdidas contrastivas, evaluado en SumMe
\cite{gygli2014summe}, TVSum \cite{song2015tvsum} y conjuntos de resumen multimodal con
salida textual.

\subsection{Análisis de resultados}
El artículo reporta F1 $\approx 55.0$ en SumMe y $\approx 63.0$ en TVSum, por encima de
métodos previos basados en dispersión o grafos \cite{chu2018spell}. El grupo replicó el
pipeline completo y obtuvo los mismos números que el checkpoint oficial (F1 0.625 en TVSum;
0.6383 en el split 0, idéntico al publicado), lo que valida el código y el protocolo de
evaluación. El modelo trata el video como una secuencia plana de frames y sus números
provienen de corridas individuales, sin varianza entre semillas, que en conjuntos de pocas
decenas de videos cambia la lectura de los resultados.

\section{Extensión A2Summ-X y el diagnóstico contra la fusión}

A2Summ-X, trabajo previo directo de esta tesis, extiende el modelo base hacia una
representación por tomas y jerárquica: adaptadores por modalidad, fusión jerárquica de frames
a tomas y de tomas a video, alineación entre modalidades, compuertas aprendidas por modalidad
y un agregador temporal de clips. Cada componente se activa por separado.

\subsection{Análisis de resultados}
Con 25 corridas por configuración, pruebas t pareadas sobre F1 y métricas de ranking de
Otani \cite{otani2016deep}, se obtuvo que: (i) el texto aporta $\Delta$F1 $= +0.033 \pm
0.062$ ($p = 0.014$) en TVSum, con tendencia positiva no significativa en SumMe; (ii)
reemplazar GoogleNet por DINOv3 o V-JEPA-2 no cambia el resultado ($|\Delta \text{F1}| <
0.001$, $p > 0.8$); y (iii) las mejoras de arquitectura no superan la varianza entre
corridas. El análisis posterior del modelo mostró por qué la señal visual pesaba tan poco: la
alineación inyecta la transcripción en la consulta visual, la compuerta softmax reparte un
peso total fijo entre modalidades (con pesos de texto medidos entre 0.75 y 0.89) y la cabeza
de decisión solo lee la mezcla resultante. Cuatro mediciones independientes coincidieron con
ese diagnóstico: los pesos de la compuerta, el colapso visual de GraphVSum en BLiSS, la
correlación de importancia humana $\rho \approx 0$ fuera de dominio y las secuencias sin
cortes detectables. La conclusión fue que el problema está en el diseño de la fusión y que
cualquier arquitectura que mezcle las representaciones antes de decidir puede reproducirlo.

\section{Segmentación auto-supervisada de escenas: BaSSL}

Mun et al.\ \cite{mun2022bassl} entrenan con tareas contrastivas sensibles a las fronteras
sobre MovieNet \cite{huang2020movienet} y reportan el punto de comparación habitual del área
(F1 47.0 y AP 57.4 en la convención de 318 películas). TranS4mer \cite{islam2023trans4mer} y
MEGA \cite{sadoughi2023mega} refinan la arquitectura temporal y la agregación multimodal;
SCRL \cite{wu2022scrl} y UBoCo \cite{kang2022uboco} exploran la consistencia de escena sin
etiquetas.

\subsection{Análisis de resultados}
Estos sistemas predicen fronteras entre escenas contiguas. No agrupan segmentos del mismo
tema separados en el tiempo ni producen una estructura navegable por tópico. MEGA, el más
cercano a CAVA en el uso de varias modalidades, agrega y destila las modalidades hacia una
representación compartida, que es el diseño cuyo problema se documentó en A2Summ-X. Ninguno
de estos trabajos reporta qué modalidades sostienen cada escena detectada.

\section{Modelos visión--lenguaje para segmentación: Scene-VLM}

Berman et al.\ \cite{berman2025scenevlm} presentan Scene-VLM, un VLM ajustado para segmentar
video en escenas. El modelo recibe frames, la transcripción del diálogo y metadatos
opcionales de tomas consecutivas, y predice las fronteras de forma secuencial con un puntaje
de confianza y una justificación en lenguaje natural. Reportan el mejor resultado publicado
en MovieNet-318 (F1 62.1 y AP 66.8).

\subsection{Análisis de resultados}
El trabajo muestra que un VLM puede emitir juicios semánticos sobre unidades temporales, a
cambio del costo de entrenar y ejecutar un modelo grande sobre cada toma, y su salida sigue
siendo una única secuencia de fronteras. CAVA toma dos decisiones a partir de esto. Primero,
usa el VLM (Qwen2.5-VL \cite{bai2025qwen25vl}) congelado y fuera del camino crítico: sus
descripciones por bloque son una señal más que vota. Segundo, no deja la frontera fina en
manos del VLM, porque el trabajo previo midió que su sensibilidad temporal es baja (invertir o
barajar los frames cambia el score solo entre 0.020 y 0.024).

\section{Capítulos de video y tópicos no contiguos: VidChapters}

Ventura et al.\ \cite{ventura2024vidchapters} construyen VidChapters-7M, un corpus de videos
con capítulos anotados por sus propios autores, y entrenan modelos que generan capítulos a
partir de la transcripción. El problema se formula como predicción de fronteras con títulos.

\subsection{Análisis de resultados}
El corpus muestra que los autores de video declaran estructura temática a gran escala. La
formulación es secuencial (una lista de capítulos contiguos) y solo usa texto: no modela
retornos de tema, no usa imagen ni voz y supone un único nivel de granularidad. Para CAVA,
VidChapters es la referencia externa en el dominio de video web, con los capítulos de autor
como verdad de terreno, y sirve de contraste con un mapa que enlaza segmentos no adyacentes,
usa varias señales y produce dos niveles.

\section{Segmentación temática de texto}

La segmentación temática se estudió primero en texto. TextTiling \cite{hearst1997texttiling}
compara bloques adyacentes por similitud léxica y coloca fronteras en los valles de esa
curva; C99 \cite{choi2000advances} ordena una matriz de similitud entre oraciones y la divide
de forma recursiva.

\subsection{Análisis de resultados}
El score de frontera de CAVA por señal sigue la misma idea que TextTiling: un desplazamiento
entre ventanas vecinas de embeddings, con fronteras en los picos. La diferencia está en que
CAVA aplica ese cálculo por separado a cada señal y no solo a la transcripción, y en que
permite que un segmento vuelva a un tópico anterior. Ambos métodos producen una partición
contigua, lo que deja fuera los retornos de tema que interesan aquí.

\section{Fusión temprana y fusión tardía}

Snoek et al.\ \cite{snoek2005early} comparan, en análisis semántico de video, la fusión
temprana (concatenar los descriptores antes de aprender) con la tardía (aprender por
modalidad y combinar las decisiones). La revisión de Baltrušaitis et al.\
\cite{baltrusaitis2019multimodal} clasifica la fusión en métodos independientes del modelo
(temprana, tardía e híbrida) y métodos basados en el modelo.

\subsection{Análisis de resultados}
En Snoek et al.\ la fusión tardía gana en la mayoría de los conceptos, aunque con
diferencias por concepto, y la revisión de Baltrušaitis et al.\ no establece un ganador
general. CAVA pertenece a la familia de fusión tardía: las señales se combinan solo en el
nivel de las decisiones (cortes y etiquetas), nunca en el de las representaciones. Lo que
agrega frente a la fusión tardía clásica es que la decisión combinada es una partición con
retornos y no una etiqueta por muestra, y que una señal puede abstenerse en los tramos donde
no tiene evidencia.

\section{Ensambles de agrupamientos y matriz de co-asociación}

Strehl y Ghosh \cite{strehl2002cluster} formulan el problema de combinar varias particiones
de los mismos datos en una partición de consenso sin acceder a los descriptores originales.
Fred y Jain \cite{fred2005evidence} proponen la acumulación de evidencia: cada partición vota
por los pares de elementos que pone juntos, los votos se acumulan en una matriz de
co-asociación y la partición final se obtiene agrupando esa matriz con enlace simple o
promedio.

\subsection{Análisis de resultados}
La matriz de votos de CAVA, $V_{ij}$, es una matriz de co-asociación ponderada en la que cada
señal aporta una partición, y su modo de enlace promedio corresponde a la variante de Fred y
Jain que evita el encadenamiento del enlace simple. CAVA adapta este marco de tres maneras:
(i) las particiones de cada señal son temporales y se calculan con su propio percentil; (ii)
una señal sin evidencia en una unidad no vota sobre los pares que la incluyen, de modo que el
consenso se calcula solo sobre los votantes activos; y (iii) cada grupo final conserva la
contribución de cada señal. En la literatura revisada, los ensambles de agrupamientos no se
han aplicado a la segmentación temática de video con señales heterogéneas.

\section{SegLoop: del bucle de destilación al motor de señales}

SegLoop, el sistema previo del grupo, validó un bucle de dos redes para segmentar escenas: un
juez VLM congelado genera pseudoetiquetas con las que se reentrena un segmentador ligero (un
transformer temporal sobre tokens de tomas). En MovieNet-318 obtuvo una base de F1 $45.22 \pm
0.63$ y AP $43.10 \pm 0.62$ y, con preentrenamiento sensible a fronteras y una cabeza
contextual, F1 $51.97 \pm 0.47$ y AP $52.69 \pm 0.43$ en test. La fusión de imagen y texto
se descartó porque no superó a la rama visual.

\subsection{Análisis de resultados}
SegLoop cumple dos papeles en esta tesis. Su resultado negativo de fusión (la familia R2 con
texto no superó a R2) reproduce en segmentación lo observado en A2Summ-X para resumen: al
fusionar, la señal adicional no aporta. Además, su infraestructura está verificada de punta a
punta (detección de cortes con TransNetV2 \cite{soucekl2020transnetv2}, reconocimiento de voz
\cite{radford2023whisper}, descriptores SigLIP \cite{zhai2023siglip} y bloqueo del split de
test). CAVA la usa como motor de extracción de señales a través del contrato
\texttt{frame\_awareness} (puntos de cambio, scores, cortes certificados y tramos con huella)
sin duplicar código.

\section{Modelos visión--lenguaje como jueces}

El uso de modelos grandes como evaluadores automáticos se ha extendido en la evaluación de
texto generado. En video, VideoJudge \cite{videojudge2025} entrena jueces multimodales de 3B y
7B parámetros para puntuar respuestas de comprensión de video.

\subsection{Análisis de resultados}
La experiencia del grupo con VLM como jueces matiza estos resultados: el VLM ordena bien
pares de alternativas (ranking contrafactual casi perfecto) pero da scores absolutos poco
fiables, su sensibilidad temporal es baja y sus veredictos heredan sesgos del entrenamiento.
De ahí salen tres reglas de diseño para CAVA: (i) el VLM genera descripciones por bloque, que
votan como cualquier otra señal, y no veredictos finales; (ii) cuando se usa un juez para
validar (G-enlace), el protocolo es ciego, con controles emparejados y pruebas exactas; y
(iii) ninguna compuerta depende de un único juez no validado.

\section{Aporte de la tesis}

A partir de esta revisión, CAVA reúne cuatro elementos que los trabajos anteriores no
combinan:
\begin{enumerate}
  \item Segmentación temática de video sin fusión de representaciones. Cada señal decide sus
        propias fronteras y etiquetas, y el tópico surge de una matriz de co-asociación con
        quórum \cite{fred2005evidence}, adaptada a particiones temporales y con abstención de
        las señales sin evidencia.
  \item Tópicos no contiguos con procedencia. Los segmentos separados en el tiempo se enlazan
        cuando un quórum de señales los agrupa, y cada grupo declara qué señales lo sostienen
        (acuerdo, soporte y coherencia dejando una unidad fuera), información que una
        representación fusionada no conserva.
  \item Versiones para el consumidor. Del mapa se derivan cinco formatos con presupuestos
        verificables, evaluados en cuatro tipos de contenido.
  \item Un inventario del estado de cada módulo (validado, acotado, parcial, no validado) por
        tipo de contenido, basado en compuertas pre-registradas.
\end{enumerate}

\chapter{MARCO TEÓRICO}

Este capítulo tiene dos partes. La primera presenta los fundamentos sobre los que se
construye CAVA: la segmentación por desplazamiento de embeddings, el consenso de particiones
mediante una matriz de co-asociación, el agrupamiento jerárquico y las métricas de
evaluación. La segunda describe el diseño del sistema por módulos: la extracción de señales,
el gestor multimodal sin fusión, el mapa de tópicos, la voz y los hablantes, y las versiones
por consumidor. La Figura~\ref{fig:mapa-mental} muestra el flujo completo: las
señales se calculan en paralelo, cada una vota y ninguna representación se combina con otra;
el mapa de tópicos sale del consenso y alimenta los formatos de consumidor.

\begin{figure}[H]
\centering
\includegraphics[width=\textwidth]{generadas/f21_flujo_cava}
\caption{Flujo de CAVA. El video se divide en unidades de 5\,s; cada señal calcula sus
propias fronteras y etiquetas dentro de la franja roja, donde ninguna representación se
combina con otra; la votación decide fronteras y tópicos con quórum y abstención; del mapa de
tópicos salen los cinco formatos de versión.}
\label{fig:mapa-mental}
\end{figure}

\section{Fundamentos}

Los conceptos de este capítulo se ilustran con un caso sintético: un documental de cuatro
minutos con seis escenas de 40\,s, cada una de 8 unidades. Las dos primeras tratan de leones
(tema A), las dos siguientes del océano (tema B), la quinta es un paisaje sin narración (tema
C) y la última vuelve a los leones, de modo que el tema A se retoma. Las señales visual, de
diálogo y del VLM se generan a partir de un prototipo por tema más ruido, y el diálogo no
tiene evidencia durante la escena sin narración. Los fotogramas son ilustraciones creadas
para este ejemplo y no provienen de ningún corpus; el procesamiento, en cambio, es el del
gestor de CAVA con los parámetros del perfil de documental.

\subsection{Segmentación por desplazamiento de embeddings}
Sea $x_1, \dots, x_n$ la secuencia de embeddings de una señal, uno por unidad temporal. El
score de frontera entre las unidades $t$ y $t+1$ compara la media de una ventana de $w$
unidades a cada lado:
\begin{equation}
  b_t = 1 - \cos\!\Big(\tfrac{1}{w}\textstyle\sum_{k=0}^{w-1} x_{t-k},\;
        \tfrac{1}{w}\textstyle\sum_{k=1}^{w} x_{t+k}\Big),
\end{equation}
con los $x_t$ normalizados a norma unitaria y $w = 2$ en el gestor. El score se divide por su
máximo en el video. Los picos de $b_t$ son candidatos a frontera, como en TextTiling
\cite{hearst1997texttiling}.
CAVA corta donde $b_t$ supera un percentil propio de cada señal y exige una separación mínima
entre cortes. Usar un percentil y no un umbral absoluto hace que la cantidad de cortes no
dependa de la escala de cada embedding.
La Figura~\ref{fig:m1} muestra el score de cada señal en el caso ilustrativo.

\begin{figure}[H]
\centering
\includegraphics[width=\textwidth]{generadas/m1_fronteras_por_senal}
\caption{Score de frontera $b_t$ de cada señal en el caso ilustrativo, con el color del tema
de cada escena. Las señales visual y del VLM marcan sus picos en los cambios de tema (1:20,
2:40 y 3:20). El diálogo solo tiene un pico claro en 1:20; durante la escena sin narración
(gris) su score es nulo, y el percentil lo obliga a proponer cortes de poca altura en 2:35 y
3:55.}
\label{fig:m1}
\end{figure}

\subsection{Consenso de particiones y matriz de co-asociación}
Dadas $S$ particiones de las mismas $n$ unidades, cada una con etiquetas $\ell^{s}_i$, la
matriz de co-asociación cuenta en cuántas particiones aparece cada par en el mismo grupo
\cite{fred2005evidence}:
\begin{equation}
  V_{ij} = \sum_{s=1}^{S} w_s\, \mathbf{1}\big[\ell^{s}_i = \ell^{s}_j\big].
\end{equation}
Con $w_s = 1$, $V_{ij}/S$ es la fracción de particiones que ponen juntos a $i$ y $j$. La
partición de consenso se obtiene agrupando la matriz de distancias $D = 1 - V/S$. Este
esquema combina decisiones y no representaciones: solo necesita las etiquetas de cada
partición, nunca los embeddings que las produjeron \cite{strehl2002cluster}.

Cuando una partición no tiene información sobre una unidad (por ejemplo, la transcripción en
un tramo sin diálogo), contarla como ``separa a $i$ de $j$'' introduce un desacuerdo falso.
CAVA define el conjunto de votantes activos en un par, $A_{ij} = \sum_s w_s\,
\mathbf{1}[\ell^s_i \neq -1 \wedge \ell^s_j \neq -1]$, donde $-1$ indica abstención, y usa
la fracción $V_{ij}/A_{ij}$ en lugar de $V_{ij}/S$ (Figura~\ref{fig:abstencion}).

\begin{figure}[H]
\centering
\includegraphics[width=\textwidth]{generadas/m3_matriz_votos}
\caption{Matriz de votos del caso ilustrativo, ordenada en el tiempo (48 unidades). Los
bloques oscuros de la diagonal son los temas, y los bloques fuera de la diagonal entre las
dos apariciones de A son el retorno. El recuadro rojo marca la escena sin narración: sin
abstención, el diálogo vota sobre esas unidades como si su silencio fuera información; con
abstención, la fracción se calcula solo sobre la señal visual y el VLM.}
\label{fig:abstencion}
\end{figure}

\subsection{Agrupamiento jerárquico y percolación}
El agrupamiento jerárquico aglomerativo une en cada paso los dos grupos más cercanos. Con
enlace simple, la distancia entre grupos es la del par más cercano; con enlace promedio, la
media de todos los pares. Si se umbraliza $V$ y se toman las componentes conexas del grafo
resultante, el resultado equivale al enlace simple: basta una cadena $i \sim j \sim k$ de
pares que superan el quórum para que $i$ y $k$ queden en el mismo grupo aunque casi ninguna
señal los agrupe. En una secuencia larga con retornos, estas cadenas se propagan hasta formar
un tópico gigante, fenómeno que aquí se llama \emph{percolación}. El enlace promedio exige
que el consenso se mantenga en promedio sobre todos los pares de los dos grupos, lo que corta
las cadenas \cite{fred2005evidence} (Figura~\ref{fig:m4}).

\begin{figure}[H]
\centering
\includegraphics[width=\textwidth]{generadas/m4_percolacion}
\caption{Percolación en una cadena de seis segmentos. Cada par vecino tiene el quórum de dos
señales, pero los extremos no tienen nada en común. Con componentes conexas la cadena forma
un solo tópico; con enlace promedio se separa en tres grupos, porque el consenso medio entre
grupos no alcanza el quórum.}
\label{fig:m4}
\end{figure}

\subsection{Métricas}
El coeficiente de correlación de Matthews \cite{matthews1975comparison} resume una matriz de
confusión binaria en un valor entre $-1$ y $1$:
\begin{equation}
  \mathrm{MCC} = \frac{TP\cdot TN - FP\cdot FN}
  {\sqrt{(TP+FP)(TP+FN)(TN+FP)(TN+FN)}}.
\end{equation}
CAVA lo aplica a la clasificación de pares de unidades (mismo tópico o no). A diferencia de la
exactitud o el F1, el MCC vale cero para un predictor que une todo o que separa todo, de modo
que no premia ni la sobresegmentación ni la percolación. El índice de Rand ajustado (ARI)
\cite{hubert1985comparing} compara dos particiones corrigiendo el acuerdo esperado por azar, y
se usa para la identidad de hablante. Para elegir el número de hablantes se usa el criterio de
información bayesiano (BIC) \cite{schwarz1978estimating} sobre un agrupamiento de Ward
\cite{ward1963hierarchical}, que en cada paso une los grupos cuya fusión menos aumenta la
varianza interna.

\section{Diseño del sistema: definiciones}

\begin{itemize}
  \item Unidad: tramo temporal fijo (5\,s en los perfiles de película y documental, 10\,s en
        reunión) sobre el que cada señal produce un embedding o un score. Es el elemento que
        se vota. Las fronteras finales se ajustan al corte visual certificado por SegLoop más
        cercano, dentro de una ventana de hasta 30\,s.
  \item Señal: fuente de evidencia que se procesa por separado hasta la decisión. El producto
        usa cuatro señales de embeddings (transcripción proyectada con bge, visual SigLIP
        \cite{zhai2023siglip} por unidad, voz con log-mel y tono, y descripciones de un VLM
        \cite{bai2025qwen25vl} por bloques de $\sim$40\,s) y, en el perfil de reunión, los
        turnos de hablante como votante de fronteras.
  \item Fusión: combinar embeddings o scores de señales distintas en una representación única
        antes de decidir. CAVA no la usa en ningún punto. La base de comparación
        Fusión de los experimentos concatena los embeddings normalizados de las
        señales y los pasa por el mismo gestor como si fueran una sola señal.
  \item Consenso de votos: cada señal emite sus propios cortes (percentil propio de su score
        de frontera) y sus propias etiquetas de tópico. Una frontera queda confirmada si
        alcanza la mayoría ponderada de las señales o un número mínimo fijo de señales;
        un par de unidades se enlaza si un quórum de señales lo
        agrupa.
  \item Tópico: grupo de unidades, posiblemente no contiguas, que un quórum de señales
        agrupa. Un retorno es la reaparición de un tópico después de una ausencia.
  \item Abstención: una señal sin evidencia en una unidad no vota sobre ella (etiqueta $-1$).
  \item Procedencia: por cada grupo, el aporte de cada señal medido como acuerdo (fracción de
        pares del grupo que la señal pone en su mismo tópico propio), soporte (fracción de
        unidades que comparten etiqueta con otra del grupo) y coherencia de sus embeddings
        dejando una unidad fuera.
  \item Versión por consumidor: producto derivado del mapa con un presupuesto declarado:
        capítulos, recapitulación de 60\,s, destacados del 10\,\%, hilos de tópico e hilos
        por hablante.
\end{itemize}

\section{Extracción de señales}

La transcripción usa faster-whisper large-v3
\cite{radford2023whisper} y se proyecta a embeddings con bge \cite{xiao2023bge}. La señal
visual codifica un frame por unidad con SigLIP. La voz extrae descriptores de log-mel y tono.
El VLM (Qwen2.5-VL-7B, congelado, sin ajuste) describe bloques de $\sim$40\,s a partir de
tres frames y del habla del bloque, y esas descripciones, proyectadas con bge, forman una
cuarta señal de embeddings. El diseño admite además otras señales con la misma interfaz;
el texto en pantalla (OCR) y el espectro de audio están definidos pero no implementados. Los perfiles de contenido fijan
el percentil de corte por dominio: p80 para película, p90 para documental y p95 para reunión.

\section{Gestor multimodal sin fusión}

El gestor, núcleo del sistema, ejecuta cuatro pasos:

\begin{enumerate}
  \item Por señal: score de frontera propio (ecuación del desplazamiento con ventana $w$),
        cortes propios por percentil con separación mínima y etiquetas de tópico propias
        sobre esos cortes.
  \item Consenso de fronteras: cada señal marca presencia en $\pm 1$ unidad alrededor de
        cada corte propio, los votos se suman con el peso de cada señal y cada racha de
        posiciones contiguas que alcanza la mayoría se reduce a un único corte en el máximo
        de votos. En el producto todas las señales de embeddings pesan 1.0; el votante de
        densidad de habla, usado en los pilotos, pesa 0.5.
  \item Consenso de enlace: la matriz de votos $V_{ij}$ se compara con el quórum $q$ y los
        grupos se forman como componentes conexas o por enlace promedio. Con abstención, un par solo se enlaza si tiene al
        menos $q$ votantes activos.
  \item Procedencia y banderas: cada grupo reporta el aporte de cada señal; un grupo de una
        unidad sin soporte de varias señales se marca como sin tópico, y un grupo con más
        segmentos que un máximo fijado se marca para revisión.
\end{enumerate}

En el caso ilustrativo (Figura~\ref{fig:m2}), las tres señales coinciden en 1:20, 2:40 y
3:20. El corte del diálogo en 2:35 cae dentro de la tolerancia del corte de 2:40 y se suma a
él, mientras que el de 3:55 queda con un solo voto y se descarta.

\begin{figure}[H]
\centering
\includegraphics[width=\textwidth]{generadas/m2_consenso_fronteras}
\caption{Consenso de fronteras en el caso ilustrativo. Arriba, los cortes propios de cada
señal y su ventana de tolerancia; abajo, los votos por posición. Solo las posiciones que
alcanzan la mayoría (dos de tres señales) producen un corte, ubicado en el máximo de votos de
cada racha.}
\label{fig:m2}
\end{figure}

\section{Retornos de tópico}

Las etiquetas de cada señal admiten retornos. En orden
temporal, cada segmento puede unirse al tópico anterior con el que tenga mayor similitud
coseno, si esa similitud supera un percentil de las similitudes entre segmentos no adyacentes
de la misma señal. Un segmento nunca vuelve al tópico del que acaba de separarse y solo puede
unirse a un tópico existente. De este modo, que un tema se retome es una decisión de cada
señal, que el quórum después confirma o descarta.
La Figura~\ref{fig:m5} muestra el resultado completo sobre el caso ilustrativo: las tres
señales reconocen el regreso a los leones, el mapa final une los dos tramos del tema A en un
solo tópico y la procedencia indica que el grupo del paisaje está sostenido solo por la señal
visual y el VLM, porque el diálogo se abstiene.

\begin{figure}[H]
\centering
\includegraphics[width=\textwidth]{generadas/m5_mapa_topicos}
\caption{Mapa de tópicos del caso ilustrativo. Las tres primeras filas son las etiquetas
propias de cada señal, coloreadas según el tema de sus unidades (rayado: abstención). La
fila \emph{mapa} es el resultado del gestor con enlace promedio y abstención, y el arco une
los dos tramos del tema A. Abajo, el soporte de cada señal en cada grupo.}
\label{fig:m5}
\end{figure}

\section{Mapa de tópicos}

El mapa de tópicos es la salida que usan los formatos. Por video tiene dos niveles de
granularidad (uno fino, con quórum de todas las señales, y uno por defecto, con quórum 2) y, por tópico, las unidades que lo forman, sus tramos en segundos, la duración
acumulada, el número de retornos, una etiqueta generada por un modelo de lenguaje
(Qwen2.5-14B) a partir de los tramos, la procedencia por señal, las banderas y el contexto de
cada tramo (la descripción del VLM y el diálogo, ambos producidos por el propio sistema). Los
nombres de personajes se propagan por los enlaces del mapa. Las identidades por apariencia
(P\#) se calculan pero no se publican como identidad.

\section{Agregador adaptativo}

Como alternativa al percentil fijo, el agregador adaptativo deja que la longitud de los
tópicos dependa de dónde las señales siguen viendo continuidad. Cada votante juzga cada par
de unidades vecinas comparando la similitud cruzada con su coherencia interna y un umbral
propio, la unión del par requiere mayoría estricta de votantes activos y cada ronda se
registra como un nivel de una jerarquía. Esta variante trata la granularidad como un
resultado y no como un hiperparámetro; no forma parte del producto congelado.

\section{Voz e identidad de hablante}

La voz entra al sistema de tres maneras. Como votante de tópico, sus embeddings de log-mel y
tono agrupan unidades acústicamente parecidas. Como votante de fronteras, la caída de la
densidad de habla marca pausas. Como identidad, los embeddings de hablante WavLM-sv
\cite{chen2022wavlm} se agrupan con Ward, el número de grupos se elige con BIC y el resultado
son las identidades V\#. Un detector de hablante único compara la dispersión de los
embeddings con un umbral $\tau$ calibrado y decide si el video tiene un solo hablante, como un
narrador. En el perfil de documental las identidades de voz se publican como grupos acústicos
y no como personas, porque el detector no reconoce al narrador único.

\section{Versiones por consumidor}

Del mapa se derivan cinco formatos con presupuestos verificables: capítulos (fronteras del
nivel superior), recapitulación de 60\,s (una unidad por tópico hasta llenar el presupuesto),
destacados del 10\,\% (tramos representativos de los tópicos más largos), hilos de tópico
(una lista de reproducción por tópico con sus retornos) e hilos por hablante (una lista por
identidad de voz). La compuerta G-versión verifica, por video y por formato, que cada lista
exista, cumpla su presupuesto y que cada perfil tenga al menos una versión. No mide la
calidad percibida, que corresponde a la evaluación con personas (EXP-040).

\chapter{MARCO METODOLÓGICO}

La metodología adapta el esquema de compuertas incrementales de los trabajos previos a un
sistema que no se entrena de punta a punta. Cada capacidad de CAVA se valida con una
compuerta pre-registrada: el criterio, los datos y el umbral se fijan antes de ejecutar, y el
split de test de cada dominio se consulta una sola vez por hipótesis. El pipeline es
reproducible: cada corrida guarda un manifiesto, los experimentos se registran en orden
cronológico (EXP-009 a EXP-040), las decisiones de arquitectura se documentan como ADR y una
suite de pruebas automatizadas cubre los módulos centrales. El Anexo~C resume los
experimentos citados.

\section{Datos y splits}

\begin{itemize}
  \item MovieNet \cite{huang2020movienet}, dominio película. Las anotaciones de escena con
        la convención de BaSSL \cite{mun2022bassl} se usan para construir una referencia
        aproximada del mapa de tópicos: dos tomas pertenecen al mismo tópico si comparten
        reparto (proxy de reparto). Los splits de desarrollo (10 películas), confirmación y
        test (47 películas) son disjuntos y el de test no se usa para ajustar nada.
  \item BBC Planet Earth: 11 episodios documentales de $\sim$49 minutos con anotaciones de
        escena de 5 anotadores. Valida el perfil de documental contra el acuerdo humano.
  \item AMI \cite{carletta2007ami}: 120 reuniones de $\sim$35 minutos con anotación de
        tópicos y de hablantes. Valida el perfil de reunión y la identidad de hablante.
  \item VidChapters \cite{ventura2024vidchapters}, dominio video web: blind40 (40 videos de
        10 a 40 minutos) para la evaluación ciega contra capítulos de autor, y raw\_ext (50
        videos) para desarrollo.
  \item TVSum \cite{song2015tvsum} y SumMe \cite{gygli2014summe}: video de usuario de 1 a 11
        minutos con importancia anotada por personas, usada como referencia automática de
        la calidad de los destacados.
  \item VoxConverse \cite{chung2020voxconverse} y AVA-Speech: conjuntos de audio para la
        identidad de hablante y los filtros de habla.
  \item Conjunto A/B de EXP-040: 9 videos educativos de YouTube (divulgación, comentario,
        conversación y tutorial; 5 en español y 4 en inglés; de 21 a 57 minutos) sin
        anotación previa, reunidos para la evaluación con personas.
\end{itemize}

Ninguna corrida del producto usa anotaciones en la inferencia: el contrato de corrida
\texttt{cava-run-v1} registra \texttt{gt\_used: false} en el manifiesto de cada una (115
corridas).

\section{Pipeline experimental}

\begin{enumerate}
  \item Unidades y señales: ffprobe y ffmpeg preparan el audio y las unidades fijas;
        faster-whisper large-v3 transcribe; SigLIP codifica un frame por unidad; se extraen
        log-mel y tono; el VLM (Qwen2.5-VL-7B) describe bloques de $\sim$40\,s.
  \item Gestor: la función \texttt{manage()} se ejecuta con los hiperparámetros del perfil
        (percentil de corte, quórum, tolerancia, modo de enlace, percentil de retorno y
        abstención) y las fronteras se ajustan a los cortes visuales certificados por
        SegLoop.
  \item Mapa y etiquetas: \texttt{topic\_map.py} produce los dos niveles, los retornos y la
        procedencia; Qwen2.5-14B etiqueta los tópicos; los nombres de personajes se propagan
        por los enlaces.
  \item Versiones: \texttt{versions.py} genera los cinco formatos con clips mp4, listas de
        reproducción e índice, y G-versión verifica los presupuestos por video.
  \item Evaluación: un script por experimento lee los JSON de \texttt{results/}, calcula
        los criterios pre-registrados y el veredicto se anota en el registro de
        experimentos. Las recompensas de clips se calculan en un archivo aparte
        (\texttt{rewards.jsonl}) para que la evaluación no modifique la corrida.
\end{enumerate}

\section{Protocolo estadístico}

\begin{itemize}
  \item Pre-registro: cada compuerta declara criterio, dirección, conjuntos y umbral antes
        de abrir los datos de evaluación, con fecha en el registro de experimentos. El
        umbral típico es $z > 2$, donde $z$ es la diferencia media pareada por video
        dividida por su error estándar.
  \item Splits: el de desarrollo se usa para diagnosticar y calibrar, el de confirmación
        para decidir y el de test solo para la hipótesis final. EXP-018 confirmó en test (47
        películas) lo diagnosticado en desarrollo y decidido en confirmación.
  \item Comparaciones fuera del pre-registro: se reportan como análisis exploratorio y no
        como validación. Una confirmación previa sobre el split de entrenamiento quedó
        marcada como post hoc.
  \item Revisiones con juez: los revisores no saben de qué brazo viene cada caso, los
        controles se emparejan por película y desfase, y la significación se calcula con la
        prueba exacta de Fisher. El primer lote de G-enlace tuvo revisores distintos por
        brazo; se identificó el problema y se lanzó un segundo lote con el protocolo
        corregido.
\end{itemize}

\section{Atribución de la ganancia del gestor}

El gestor de CAVA difiere de la base \texttt{fused} en dos aspectos a la vez: no combina las
señales y, en su configuración final, usa enlace promedio con abstención. La comparación
pre-registrada en EXP-018 (\texttt{cava\_abst} frente a \texttt{fused}) mide el efecto
conjunto. Para separar los dos aspectos se reportan también, sobre las mismas 47 películas,
las comparaciones de las variantes sin abstención (\texttt{cava}, con componentes conexas, y
\texttt{cava\_avg}, con enlace promedio) frente a \texttt{fused}, y de \texttt{cava\_abst}
frente a \texttt{cava\_avg}. Estas comparaciones se calcularon con el mismo script que la
principal y se presentan como análisis secundario. La abstención no tiene equivalente en la
base fusionada, porque al concatenar los embeddings deja de existir una señal individual que
pueda abstenerse.

\section{Sistema de compuertas}

\begin{table}[H]
    \centering
    \caption{Compuertas de validación de CAVA y su estado al congelar el sistema.}
    \label{tab:gates}
    \begin{tabular}{p{2.6cm}p{6.6cm}p{3.6cm}}
        \hline
        \textbf{Compuerta} & \textbf{Criterio} & \textbf{Estado} \\
        \hline
        Test del gestor & MCC(\texttt{cava\_abst}) $-$ MCC(\texttt{fused}) $> 0$ con
                      $z > 2$ en las 47 películas de test & Cumplida ($z = 2.46$) \\
        G-coherencia & Margen intra$-$inter positivo en cada señal y en cada película
                      & Cumplida (47/47 películas en las tres señales) \\
        G-versión & Toda lista existe, cumple su presupuesto y hay al menos una versión
                      por perfil & Cumplida (58/58, 80/80, 90/90, 25/25, 9/9) \\
        G-enlace-2 & Enlaces no contiguos juzgados como el mismo hilo más que momentos
                      lejanos al azar (revisión ciega, Fisher $p < 0.05$) & No cumplida
                      ($p = 0.72$) \\
        Calidad de clips & Destacados y recapitulación superan a bases simples en
                      referencias automáticas & No cumplida (EXP-031/037/038) \\
        Narrador (EXP-039) & Un detector con habla filtrada detecta al hablante único
                      mejor que el actual ($z > 2$) sin más falsos & No cumplida
                      ($z = -0.37$ y $0.00$) \\
        EXP-040 (A/B) & Preferencia humana de las versiones frente a bases propias &
                      Diseñada; pre-registro pendiente \\
        \hline
    \end{tabular}
\end{table}

\section{Infraestructura y reproducibilidad}

Los experimentos se ejecutan en el clúster institucional (GPU NVIDIA RTX A6000 de 48\,GB y
A100 vía SLURM) con caché local de HuggingFace. El VLM y los modelos de voz corren en GPU y el
gestor en CPU. Cada corrida del producto escribe un manifiesto (commit de git, indicador de
cambios sin confirmar, perfil e hiperparámetros) y el archivo \texttt{clips.jsonl}. La suite
de 25 pruebas automatizadas cubre el gestor, el agregador adaptativo, el ajuste a cortes, el
mapa y la compuerta de versiones, y se ejecuta antes de cada envío a la cola. El paquete
\texttt{topicagg} declara sus dependencias en \texttt{pyproject.toml}. La decisión de no
fusionar representaciones y el uso de votantes activos con abstención están documentados en
la ADR-0001.

\section{Métricas}

\begin{itemize}
  \item Mapa de tópicos: MCC sobre la clasificación de pares de unidades (mismo tópico o
        no) contra la referencia (proxy de reparto en MovieNet; anotación de tópicos en
        AMI). El número de tópicos y la tasa de enlace (fracción de pares enlazados) se usan
        para diagnosticar sobresegmentación y percolación. La precisión de enlace se compara
        con la tasa base de pares positivos.
  \item Coherencia: margen entre la similitud coseno dentro de cada grupo (dejando una unidad
        fuera) y la similitud entre grupos, por señal, con la fracción de películas en que es
        positivo.
  \item Identidad de hablante: ARI contra la anotación (AMI, VoxConverse) y especificidad de
        las identidades publicadas.
  \item Fronteras: F1 contra cortes de referencia (VidChapters, BBC) y distancia del ajuste a
        los cortes visuales certificados.
  \item Calidad de clips: recompensa igual al $z$ de la fracción de anotadores que eligió cada
        segundo (TVSum, SumMe), comparada con bases de clips equiespaciados, clip central y
        azar.
  \item Versiones: existencia y cumplimiento del presupuesto de cada lista (G-versión), por
        video y formato.
\end{itemize}

\chapter{RESULTADOS}

Este capítulo reporta el estado experimental de CAVA al congelar el sistema (commit
\texttt{d51d201}, 25 pruebas automatizadas en verde). Los números provienen de corridas
completas guardadas en \texttt{results/} y del registro de experimentos EXP-009 a EXP-039.

\section{Experimentos y resultados}

\subsection{Pilotos del gestor sin fusión}

El mecanismo de consenso se verificó primero con datos sintéticos de dos señales y después en
video real. En 3 videos de TVSum, con transcripción, señal visual SigLIP/BaSSL y densidad de
voz, los grupos resultantes siguieron los pasos del contenido (por ejemplo, seguridad,
evaluación del clavo, inflado, solución de problemas y taller en un video sobre reparar un
neumático) y cada grupo tuvo soporte de varias señales. El piloto en BLiSS, con
transcripciones reales, mostró el primer problema de calibración: un embedding de texto
(RoBERTa) que se agrupa por locutor fragmenta el contenido en exceso (17 tópicos por clip), lo
que motivó los percentiles por perfil. En una validación paralela fuera de dominio, las
fronteras fueron estables en 4 dominios (AP de 0.73 a 0.78) y la correlación con la
importancia humana fue $\rho \approx 0$, lo que indica que una frontera temática y un momento
importante son cosas distintas.

\subsection{Mapa de tópicos en test (EXP-017 y EXP-018)}

EXP-017 encontró en desarrollo que el enlace por componentes conexas con retornos percola: sin
abstención, CAVA colapsa en unos 4 tópicos gigantes con una tasa de enlace de 0.99. EXP-018
pre-registró la corrección (abstención de las señales sin evidencia y enlace promedio) y la
evaluó en el split de test de 47 películas de MovieNet (Tabla~\ref{tab:exp018}).

\begin{table}[H]
    \centering
    \caption{Mapa de tópicos en test (47 películas, EXP-018). MCC por pares de unidades
    contra el proxy de reparto, número medio de tópicos por película y tasa de enlace.}
    \label{tab:exp018}
    \small
    \begin{tabular}{lccc}
        \hline
        \textbf{Sistema} & \textbf{MCC} & \textbf{N.$^{\circ}$ tópicos} & \textbf{Tasa de enlace} \\
        \hline
        \texttt{cava\_abst} (enlace promedio y abstención) & \textbf{0.0113} & 139.4 & 0.044 \\
        \texttt{cava\_abst\_r95} (percentil de retorno 95) & 0.0115 & 216.7 & 0.020 \\
        \texttt{cava} (componentes, sin abstención) & 0.0090 & 4.0 & 0.990 \\
        \texttt{cava\_avg} (enlace promedio, sin abstención) & 0.0083 & 90.7 & 0.054 \\
        visual sola & \underline{0.0195} & 56.6 & 0.248 \\
        VLM sola & 0.0079 & 32.4 & 0.159 \\
        \texttt{fused} (concatenación de señales) & 0.0053 & 39.7 & 0.197 \\
        diálogo sola & 0.0053 & 81.4 & 0.331 \\
        \texttt{contiguous} (gestor sin retornos) & 0.0045 & 86.0 & 0.001 \\
        \hline
    \end{tabular}
\end{table}

El criterio pre-registrado se cumplió: MCC(\texttt{cava\_abst}) $-$ MCC(\texttt{fused})
$= +0.0060 \pm 0.0024$ ($z = 2.46$), y frente al gestor contiguo la diferencia es $+0.0068
\pm 0.0027$ ($z = 2.48$). La Tabla~\ref{tab:ablacion} y la Figura~\ref{fig:ablacion}
separan las dos diferencias entre \texttt{cava\_abst} y \texttt{fused}. Sin abstención, no
fusionar las señales no basta para superar a la fusión: la diferencia es de $z = 0.80$ con
componentes conexas y de $z = 1.11$ con enlace promedio. La abstención agrega $+0.0030$
sobre el enlace promedio ($z = 1.78$), y solo la combinación de las dos cruza el umbral. Como
la abstención requiere que cada señal conserve su propia partición, la ganancia medida sí
depende de no fusionar, aunque por una vía indirecta.

\begin{table}[H]
    \centering
    \caption{Comparaciones pareadas en las 47 películas de test (EXP-018). Diferencia media de
    MCC por película $\pm$ error estándar. Solo la primera fila es el criterio
    pre-registrado; el resto es análisis secundario con el mismo script.}
    \label{tab:ablacion}
    \small
    \begin{tabular}{lcc}
        \hline
        \textbf{Comparación} & \textbf{Diferencia de MCC} & $z$ \\
        \hline
        \texttt{cava\_abst} $-$ \texttt{fused} (pre-registrada) & $+0.0060 \pm 0.0024$ & 2.46 \\
        \texttt{cava\_abst} $-$ \texttt{contiguous} & $+0.0068 \pm 0.0027$ & 2.48 \\
        \texttt{cava} $-$ \texttt{fused} & $+0.0037 \pm 0.0046$ & 0.80 \\
        \texttt{cava\_avg} $-$ \texttt{fused} & $+0.0031 \pm 0.0028$ & 1.11 \\
        \texttt{cava\_abst} $-$ \texttt{cava\_avg} & $+0.0030 \pm 0.0017$ & 1.78 \\
        \texttt{cava\_abst} $-$ visual sola & $-0.0082 \pm 0.0030$ & $-2.74$ \\
        \hline
    \end{tabular}
\end{table}

\begin{figure}[H]
\begin{center}
\begin{tikzpicture}
\begin{axis}[
  xbar, width=0.82\textwidth, height=6.2cm,
  xlabel={Diferencia de MCC por película (barra: $\pm 1.96$ EE)},
  symbolic y coords={absvis,abscava,avgfus,cavafus,abscont,absfus},
  ytick=data,
  yticklabels={\texttt{cava\_abst} $-$ visual, \texttt{cava\_abst} $-$ \texttt{cava\_avg},
               \texttt{cava\_avg} $-$ \texttt{fused}, \texttt{cava} $-$ \texttt{fused},
               \texttt{cava\_abst} $-$ \texttt{contiguous}, \texttt{cava\_abst} $-$ \texttt{fused}},
  yticklabel style={font=\footnotesize},
  xmin=-0.016, xmax=0.016,
  xtick={-0.015,-0.01,-0.005,0,0.005,0.01,0.015},
  xticklabel style={font=\footnotesize, /pgf/number format/fixed,
                     /pgf/number format/precision=3},
  scaled x ticks=false,
  bar width=7pt, enlarge y limits=0.12,
  extra x ticks={0}, extra x tick labels={},
  extra x tick style={grid=major, major grid style={black!50}},
]
\addplot[fill=cyan!35, draw=black!60,
  error bars/.cd, x dir=both, x explicit]
  coordinates {
    (-0.0082,absvis)  +- (0.0059,0)
    (0.0030,abscava)  +- (0.0033,0)
    (0.0031,avgfus)   +- (0.0055,0)
    (0.0037,cavafus)  +- (0.0090,0)
    (0.0068,abscont)  +- (0.0053,0)
    (0.0060,absfus)   +- (0.0047,0)
  };
\end{axis}
\end{tikzpicture}
\caption{Diferencias pareadas de MCC en el test de MovieNet (47 películas). Solo las
comparaciones de \texttt{cava\_abst} con \texttt{fused} y con \texttt{contiguous} quedan por
encima de cero; la señal visual sola supera a \texttt{cava\_abst}.}
\label{fig:ablacion}
\end{center}
\end{figure}

Los valores absolutos de MCC son bajos (alrededor de 0.01) porque el proxy de reparto mide
coaparición de personajes y no tópicos, y porque la tasa de pares positivos es alta. La señal
visual sola obtiene un MCC mayor que el consenso ($z = -2.74$ para la diferencia), lo que
indica que, en este proxy, las otras señales restan más de lo que aportan. En precisión de
enlace frente a la tasa base de pares positivos, \texttt{cava\_abst} supera a la base por
$+0.026$ ($z = 5.59$), la visual sola por $+0.016$ ($z = 4.31$) y \texttt{fused} por $+0.005$
($z = 1.95$). Lo que el consenso agrega sobre la señal visual sola es estructural: la
abstención elimina la percolación (de 4 a 139 tópicos y de 0.99 a 0.04 de tasa de enlace) y
cada grupo declara qué señales lo sostienen.

La compuerta G-coherencia se cumplió en las 47 películas: el margen intra$-$inter de
\texttt{cava\_abst} es positivo en todas para las tres señales, con medias de $+0.053$
(visual), $+0.018$ (diálogo) y $+0.121$ (VLM), frente a $+0.027$, $+0.007$ y $+0.067$ de la
fusión. El gestor contiguo tiene márgenes mayores ($+0.067$, $+0.049$ y $+0.155$), lo que es
esperable porque sus grupos están formados solo por unidades vecinas, que se parecen más
entre sí.

\subsection{Perfiles de contenido y ajuste de fronteras}

El perfil de documental se validó en BBC Planet Earth contra el acuerdo de 5 anotadores ($z =
6.1$ y $7.4$ en los dos criterios), con una estructura que alcanza el 62\,\% del acuerdo
humano. El ajuste de fronteras a cortes visuales certificados se validó en 91 películas ($z
= 26.6$ y $23.0$): las fronteras consensuadas caen en cortes reales sin desplazamiento
sistemático. En reuniones (AMI), el perfil de reunión iguala a la partición contigua
($+0.036 \pm 0.035$, $z = 1.0$) y los turnos de hablante como votante de fronteras aportan
$+0.075$ ($z = 2.32$). En este dominio la estructura temática es casi contigua, y el sistema
devuelve pocos enlaces no contiguos.

\subsection{Voz: votante, identidad y hablante único}

Como votante de tópico, la voz distingue pares de unidades con el mismo hablante con un AUC
de 0.72 a 0.74, frente a 0.55 de la señal visual (EXP-020a). La identidad de hablante
(WavLM-sv, Ward y BIC) supera a la agrupación por apariencia: ARI 0.186 frente a 0.038 en
desarrollo de AMI, con una especificidad de las identidades publicadas del 85\,\% (69 de 81,
$p \approx 0$, binomial). Con los filtros de habla entrenados, la identidad de voz llega a ARI
0.64 en VoxConverse y 0.73 en AMI (desde 0.25 y 0.29). El detector de hablante único no da
falsos positivos en 40 reuniones de AMI, pero tampoco detecta al narrador en ninguno de los
11 episodios de BBC; por eso el perfil de documental publica grupos acústicos y no personas.
La compuerta de AMI con el mapa completo resultó no inferior a la configuración base
($-0.007 \pm 0.009$, dentro del margen).

EXP-039 probó cinco detectores alternativos del hablante único sobre habla filtrada. En
desarrollo se eligió el de dispersión (\emph{spread}) con $\tau = 0.983$. En confirmación,
el detector no mejora la detección del hablante único frente al producto actual, ni con el
filtro S\&W ($-0.022 \pm 0.058$, $z = -0.37$) ni con I\&W ($0.000 \pm 0.062$, $z = 0.00$),
aunque cumple los criterios de falsos positivos (a lo sumo 0.019 en mezclas con varias voces
y ninguna reunión de AMI marcada). Según la lectura declarada en el pre-registro, ninguno de
los detectores probados conserva la detección del narrador, y el módulo queda como no
validado.

\subsection{Compuerta de versiones (G-versión)}

La verificación automática de presupuestos se cumplió en todos los conjuntos: 58 de 58
películas de MovieNet, 80 de 80 reuniones de AMI, 90 de 90 corridas del contrato v1 (los 40
videos web de VidChapters blind40 y los 50 de TVSum), 25 de 25 videos de SumMe y los 9 videos
del conjunto A/B. G-versión verifica existencia y presupuesto, no calidad percibida.

\subsection{Enlaces no contiguos (G-enlace)}

El primer lote de revisión con juez modelo dio 83\,\% frente a 63\,\% (Fisher $p = 0.072$),
con pocos casos y con revisores distintos en cada brazo. El segundo lote, pre-registrado con
el protocolo corregido (ciego, controles emparejados y 60 casos nuevos por brazo), dio 62\,\%
frente a 65\,\% ($p = 0.72$), de modo que G-enlace-2 no se cumple. La revisión de los casos
indica que los enlaces de CAVA capturan continuidad de personajes y situaciones, que el juez
reconoce, pero no la continuidad narrativa que exige un juez de trama. Por eso los retornos se
presentan como reagrupación temática con procedencia y no como tramas.

\subsection{Calidad automática de clips (EXP-031, EXP-037 y EXP-038)}

Ninguna regla de selección de clips superó a las bases simples con referencias automáticas.
En TVSum (EXP-031), los destacados de CAVA no difieren del azar ($+0.017 \pm 0.076$, $z =
0.22$) y quedan por debajo de un clip único en el centro del video ($-0.265 \pm 0.117$, $z =
-2.26$). Esa base aprovecha el sesgo de posición de los anotadores de TVSum, que puntúan más
alto la parte central del video. En SumMe (EXP-037), la mejor regla (H3, ventana saliente por
tópico) obtiene 0.278 frente a 0.156 de los clips equiespaciados, sin llegar a 2 errores
estándar ($z = 0.70$). La política de clips aprendida con un bandido contextual sobre TVSum
(EXP-038) obtiene en SumMe una recompensa de $-0.265$, peor que el azar ($z = -2.62$). TVSum
premia la mitad del video y SumMe el principio, y la política aprendió preferencias de
posición propias del conjunto de entrenamiento. En SumMe apareció además un defecto del
producto: en videos de menos de 50\,s, el 10\,\% de la duración no alcanza a cubrir una unidad
de 5\,s y los destacados salen vacíos. La corrección propuesta es exigir al menos una unidad.
Con estas referencias, la importancia humana de TVSum y SumMe y los capítulos de autor, no se
distingue a CAVA de las bases simples, porque cada referencia tiene su propio sesgo de
posición.

\subsection{Video web (EXP-032) y nivel de capítulo (EXP-033)}

Contra los capítulos de autor de VidChapters blind40, CAVA no supera a los cortes uniformes
con el mismo número de cortes ($-0.009$, $z = -0.44$) ni al mejor contiguo ($+0.046$, $z =
1.1$), y sí supera a la fusión en el tercer criterio pre-registrado ($+0.067 \pm 0.027$, $z =
2.46$). El nivel de tópico resulta unas 2.7 veces más fino que los capítulos de autor. El
nivel de capítulo por fronteras (las $N$ más fuertes) cumple el criterio de F1 de fronteras
($z = 2.57$) pero no el de agrupamiento ($z = 0.59$), por lo que queda como parcial y no se
incorpora al producto.

\subsection{Reconocimiento de voz (EXP-030)}

Se evaluó Parakeet TDT 0.6B v3 \cite{koluguri2024parakeet} como reemplazo de Whisper
large-v3. Empata en habla falsa ($+0.003$, $z = 0.11$), pierde menos en el número efectivo de
hablantes ($-0.039$) y es no inferior en MCC de AMI ($+0.015$, $z = 0.45$), pero omite seis
veces más habla (6.0\,\% frente a 1.5\,\%, fuera del margen de $+0.02$). Whisper large-v3 se
mantiene como transcriptor del producto.

\subsection{Evaluación con personas (EXP-040): estado}

El diseño compara, formato por formato, las versiones de CAVA con dos bases propias que
comparten unidades, corte, codificación, presupuesto y modelo de etiquetas: U, con clips
equiespaciados sin estructura, y S, con el mapa de la mejor señal sola (visual). Se propone
una tercera base, T: la segmentación que Claude Opus 5.5 produce leyendo solo los subtítulos
de YouTube, que permite medir si CAVA aporta algo sobre un modelo de lenguaje que solo usa el
texto. Cada participante ve dos versiones del mismo formato y video, en orden aleatorio y
sin saber su origen, elige una (o empate) y puntúa la utilidad de cada una. El resultado
principal por formato es la tasa de victoria de CAVA frente a U, con intervalos por bootstrap
sobre videos y participantes, y el ranking se estima con un modelo de Bradley--Terry
\cite{bradley1952rank}. Para detectar una preferencia de 60/40 con potencia 0.8 y $\alpha =
0.05$ se necesitan unos 200 juicios por formato, cerca de 1000 en total.

Al cierre de este documento, el conjunto de 9 videos está descargado y verificado, CAVA se
ejecutó sobre los 9 con el sistema congelado (G-versión 9/9, sin cambios locales en el
repositorio), la base T está generada para los 9 videos y la página de etiquetado ciego está
lista. Quedan pendientes el pre-registro, el etiquetado propio, el piloto y la evaluación
completa.

\subsection{Inventario del sistema}

La Tabla~\ref{tab:inventario} resume el estado de los 26 módulos del pipeline al congelarlo.
El detalle de cada fila, con la evidencia que la respalda, está en \texttt{docs/progress.md}.

\begin{table}[H]
    \centering
    \caption{Estado de validación de los 26 módulos de CAVA.}
    \label{tab:inventario}
    \begin{tabular}{p{2.6cm}cp{9.4cm}}
        \hline
        \textbf{Estado} & \textbf{N.$^{\circ}$} & \textbf{Módulos} \\
        \hline
        Validado & 10 & señal visual, voz como votante de tópico, gestor sin fusión, ajuste a
                         cortes visuales, perfiles de contenido, mapa con retornos y
                         procedencia, etiquetas de tópico, nombres de personajes, identidad de
                         voz, versiones con G-versión \\
        Acotado & 2 & VLM (sirve para enlazar, no como juez de trama); detector de hablante
                      único (no detecta al narrador) \\
        Parcial & 3 & turnos de hablante en reunión, filtro de habla, nivel de capítulo \\
        No validado & 6 & pesos de voto aprendidos, identidades por apariencia, regla de
                          destacados, regla H3, política de clips aprendida, detectores de
                          segunda voz (EXP-039) \\
        En producto sin compuerta propia & 4 & unidades fijas, transcripción, clips y
                          listas de reproducción, contrato de corrida \\
        Sin evaluación & 1 & exportación de datos de entrenamiento \\
        \hline
    \end{tabular}
\end{table}

Los módulos no validados quedan fuera del producto congelado, que usa pesos iguales, no
incluye el nivel de capítulo ni el filtro de habla, y selecciona destacados con la regla de
tópicos más largos.

\section{Discusión}

Los resultados apoyan la hipótesis estructural con una salvedad importante. El consenso con
abstención produce mapas de tópicos con mejor MCC que la fusión y que el gestor contiguo en
el test pre-registrado ($z = 2.46$ y $2.48$), y conserva márgenes de coherencia positivos en
todas las películas. Sin embargo, el efecto absoluto es pequeño, no fusionar las señales sin
abstención no basta para superar a la fusión, y la señal visual sola puntúa más alto en el
proxy disponible. El valor del sistema multimodal está en propiedades que el MCC no recoge:
la abstención evita la percolación (sin ella el mapa colapsa en 4 tópicos), cada grupo
declara qué señales lo sostienen, un mismo conjunto de perfiles cubre cuatro tipos de
contenido y del mismo mapa salen cinco formatos.

Las referencias automáticas de calidad de clips no distinguen a CAVA de bases simples, y una
política entrenada en un conjunto empeora en el otro. Esto no afecta al mapa, pero obliga a
medir la calidad de las versiones de otra forma: con preferencia humana directa, formato por
formato, frente a bases propias con el mismo presupuesto. Ese es el diseño de EXP-040, cuyo
conjunto de videos ya está procesado.

Hay cuatro amenazas a la validez. El proxy de reparto de MovieNet mide coaparición de
personajes y no trama, como también mostró G-enlace-2. La base \texttt{fused} es una
concatenación sin aprendizaje, más simple que la fusión con compuertas de A2Summ-X que motivó
el trabajo, de modo que la comparación no cubre las fusiones aprendidas. La evaluación en AMI
depende de que el corpus tenga disponible el video de cada reunión, y un trabajo de
extracción falló por archivos ausentes. Por último, los dominios sin anotación de tópicos
(YouTube largo, video de usuario) solo admiten referencias sustitutas, que este trabajo trata
como no concluyentes.

\customchapter{CONCLUSIONES}

\textit{Problema de partida.} La tesis parte del problema identificado en A2Summ-X: al
fusionar señales multimodales en una representación única, unas reescriben a otras antes de
la decisión. En ese modelo el texto llegó a pesos de 0.75 a 0.89 en la compuerta y anuló la
contribución visual, y la familia de fusión de SegLoop mostró el mismo comportamiento en
segmentación de escenas. De ahí sale la regla de diseño de este trabajo: mantener separada la
identidad de cada señal hasta la decisión.

\textit{Aporte y validación.} Se construyó un sistema de segmentación temática de video largo
en el que cada señal (transcripción, imagen, voz y descripciones de un VLM) emite sus propias
fronteras y etiquetas, y el tópico surge de una matriz de co-asociación con quórum,
abstención y procedencia por grupo. En el test pre-registrado de 47 películas, el gestor
superó a la fusión por concatenación ($z = 2.46$) y al gestor contiguo ($z = 2.48$), y su
margen de coherencia fue positivo en las tres señales y en todas las películas. El análisis
secundario muestra que la ganancia sobre la fusión aparece solo cuando se combinan el enlace
promedio y la abstención, y que la señal visual sola obtiene un MCC más alto. El ajuste a
cortes visuales se validó en 91 películas, el perfil de documental se validó en BBC frente
al acuerdo humano, la identidad de hablante alcanzó un ARI de 0.73 en AMI y las cinco
versiones cumplieron su presupuesto en todos los tipos de contenido evaluados (película,
reunión, video web y video de usuario).

\textit{Estado de la evaluación.} De los 26 módulos del sistema congelado, 10 están
validados, 2 acotados, 3 parciales y 6 no validados; 4 forman parte del producto sin
compuerta propia y 1 no tiene evaluación. Tres compuertas no se cumplieron. Los enlaces no
contiguos no pasan el juez de trama (G-enlace-2, $p = 0.72$) y se publican como continuidad
de personajes y situaciones. Ninguna regla de selección de clips supera a las bases simples
con referencias automáticas (EXP-031 y EXP-037), y la política aprendida empeora al cambiar
de conjunto (EXP-038). Ningún detector probado resuelve al narrador de los documentales
(EXP-039). Además, Parakeet no reemplaza a Whisper como transcriptor (EXP-030).

\textit{Implicancias.} En los dos trabajos previos, agregar componentes sobre
representaciones fusionadas no mejoró los resultados; en este, la mejora apareció al mover
el punto en que se toma la decisión. El consenso tardío produce una estructura que se puede auditar (procedencia
por señal, abstención que evita la percolación y perfiles comunes a varios dominios), a
costa de no maximizar el proxy de una sola señal. Lo que las referencias automáticas no
pueden medir es si las personas prefieren las versiones de CAVA. Esa medición corresponde a
EXP-040, cuyo conjunto de 9 videos ya está procesado y cuyo pre-registro está pendiente.

\chapter*{\center \Large TRABAJOS FUTUROS}
\addcontentsline{toc}{section}{\bfseries TRABAJOS FUTUROS}
\markboth{TRABAJOS FUTUROS}{TRABAJOS FUTUROS}

A partir de los resultados actuales quedan definidas las siguientes líneas de trabajo:

\begin{itemize}
  \item Evaluación A/B con personas (EXP-040). Pre-registrar y ejecutar el diseño de
        $\sim$1000 juicios (200 por formato) de los cinco formatos frente a las bases U
        (clips equiespaciados con igual presupuesto) y S (mejor señal sola), y decidir si se
        incluye la base T (segmentación de un modelo de lenguaje sobre los subtítulos). Antes
        de medir hay que corregir y declarar el defecto de los destacados en videos de menos
        de 50\,s.
  \item Fusión aprendida como base de comparación. Comparar el gestor con una fusión con
        compuertas aprendidas, como la de A2Summ-X, y no solo con la concatenación, para
        saber si la ventaja se mantiene frente a fusiones más fuertes.
  \item Señales pendientes de extracción. Implementar el texto en pantalla (OCR) y el
        espectro de audio como votantes, cuyos contratos ya están definidos en
        \texttt{signals/}.
  \item Enlaces de trama (G-enlace). Probar descriptores narrativos que capturen continuidad
        de trama y no solo de personajes y situaciones, o reformular la compuerta con jueces
        humanos especializados.
  \item Narrador en documentales. Después de EXP-039, el siguiente candidato es un modelo de
        segmentación con conteo local de hablantes, como pyannote, o una estimación del
        número de hablantes por \emph{eigengap} \cite{park2020nme} sobre ventanas más
        cortas; la anotación propia de los 176 clips de BBC (EXP-035b) sigue pendiente.
  \item Recompensa para políticas de clips. Normalizar la recompensa por posición o
        entrenar con un conjunto grande, dado que la importancia anotada en un conjunto
        pequeño induce preferencias de posición que no transfieren (EXP-038).
  \item AMI y otros dominios. Completar la evaluación de AMI que quedó pendiente por archivos
        de video ausentes y extender el A/B a otros tipos de contenido sin anotación
        pública.
\end{itemize}
 % sección opcional

%% ============================================================================
\renewcommand{\bibname}{\hfill\Large\bf{REFERENCIAS BIBLIOGRÁFICAS}\hfill}

\bibliographystyle{IEEEtran} % Estableciendo el estilo de citas IEEEtram.

\bibliography{referencias} % Recibe las referencias de IEEE

\chapter*{\center \Large ANEXOS}
\addcontentsline{toc}{section}{\bfseries ANEXOS}
\markboth{ANEXOS}{ANEXOS}

\section*{Anexo A. Reproducibilidad experimental}

El repositorio del proyecto contiene el paquete \texttt{topicagg} con sus dependencias
declaradas en \texttt{pyproject.toml} y una suite de 25 pruebas automatizadas; el contrato de
corrida \texttt{cava-run-v1}, con un manifiesto por video (commit de git, indicador de
cambios sin confirmar, perfil, hiperparámetros y \texttt{gt\_used: false}) y 115 corridas
registradas; el registro cronológico de experimentos con sus pre-registros y veredictos
(\texttt{docs/EXPERIMENT\_LOG.md}); la decisión de arquitectura ADR-0001 (sin fusión de
representaciones, votantes activos y abstención); los reportes semanales en
\texttt{reports/}; y los JSON de resultados por experimento en \texttt{results/}.

\section*{Anexo B. Inventario de módulos}

El inventario de los 26 módulos del pipeline (implementación, evidencia y estado), la tabla de
lo que puede afirmarse por tipo de contenido y la lista de conjuntos con anotación disponible
están en \texttt{docs/progress.md}. La Tabla~\ref{tab:inventario} del Capítulo~IV resume ese
inventario.

\section*{Anexo C. Glosario y experimentos citados}

\begin{table}[H]
    \centering
    \caption{Términos y nombres internos usados en el documento.}
    \label{tab:glosario}
    \small
    \begin{tabular}{p{3.2cm}p{10.2cm}}
        \hline
        \textbf{Término} & \textbf{Significado} \\
        \hline
        \texttt{cava\_abst} & Configuración final del gestor: retornos con percentil 90,
                              enlace promedio y abstención. \\
        \texttt{cava}, \texttt{cava\_avg} & Variantes sin abstención, con componentes
                              conexas o con enlace promedio. \\
        \texttt{\_r80}, \texttt{\_r95} & Sufijos que cambian el percentil de retorno a 80 o 95. \\
        \texttt{fused} & Base que concatena los embeddings normalizados de las señales y los
                         trata como una sola señal en el mismo gestor. \\
        \texttt{contiguous} & Gestor sin retornos: solo produce tópicos contiguos. \\
        Proxy de reparto & Referencia de MovieNet: dos tomas son del mismo tópico si
                           comparten personajes. \\
        Procedencia & Aporte de cada señal a un grupo (acuerdo, soporte y coherencia). \\
        Percolación & Colapso del mapa en pocos tópicos gigantes por encadenamiento de
                      enlaces. \\
        P\#, V\# & Identidades por apariencia (no publicadas) y por voz. \\
        blind40, raw\_ext & Subconjuntos de VidChapters para evaluación ciega (40 videos) y
                            desarrollo (50 videos). \\
        S\&W, I\&W & Filtros de habla de EXP-034: intersección de Silero VAD (S) o de ina (I)
                        con la detección de habla de Whisper (W). \\
        Bases U, S, T & Bases de EXP-040: clips equiespaciados, mejor señal sola y
                        segmentación de un modelo de lenguaje sobre los subtítulos. \\
        \hline
    \end{tabular}
\end{table}

\begin{table}[H]
    \centering
    \caption{Experimentos citados en el documento.}
    \label{tab:experimentos}
    \small
    \begin{tabular}{p{2.3cm}p{7.4cm}p{3.6cm}}
        \hline
        \textbf{Experimento} & \textbf{Pregunta} & \textbf{Resultado} \\
        \hline
        EXP-017/018 & ¿El mapa por consenso supera a la fusión y al contiguo en MCC? &
                      Sí con abstención ($z = 2.46$) \\
        EXP-020a & ¿La voz identifica al hablante mejor que la imagen? & Sí (AUC 0.74 frente
                   a 0.55) \\
        EXP-023 & ¿Los turnos de hablante mejoran las fronteras en reuniones? & Parcial ($z =
                  2.32$; empata con contiguo) \\
        EXP-025 & ¿El perfil de documental se acerca al acuerdo humano en BBC? & Sí ($z =
                  6.1$ y $7.4$) \\
        EXP-027 & ¿El ajuste a cortes visuales es preciso? & Sí ($z = 26.6$ y $23.0$) \\
        EXP-028 & ¿Los filtros de habla mejoran la identidad de voz? & Sí (ARI 0.64 y
                  0.73) \\
        EXP-029 & ¿El detector de hablante único evita falsos en reuniones? & Sí (0/40), sin
                  narrador (0/11) \\
        EXP-030 & ¿Parakeet reemplaza a Whisper? & No \\
        EXP-031 & ¿Los destacados siguen la importancia humana en TVSum? & No \\
        EXP-032 & ¿CAVA reproduce los capítulos de autor? & Solo supera a la fusión \\
        EXP-033 & ¿Un nivel de capítulo por fronteras funciona? & Parcial \\
        EXP-037 & ¿Una regla fija de destacados supera a las bases? & No ($z = 0.70$) \\
        EXP-038 & ¿Una política aprendida transfiere entre conjuntos? & No ($z = -2.62$) \\
        EXP-039 & ¿Otro detector resuelve al narrador con habla filtrada? & No ($z = -0.37$) \\
        G-enlace-2 & ¿Un juez reconoce los retornos como el mismo hilo? & No ($p = 0.72$) \\
        EXP-040 & ¿Las personas prefieren las versiones de CAVA? & Diseñado, pendiente \\
        \hline
    \end{tabular}
\end{table}

\section*{Anexo D. Definiciones operativas de métricas}

El MCC del mapa se calcula sobre la clasificación de pares de unidades (mismo tópico o no)
contra la referencia: proxy de reparto en MovieNet y anotación de tópicos en AMI. El margen
de coherencia intra$-$inter compara, por señal, la similitud coseno de cada unidad con su
grupo (dejándola fuera) y la similitud entre grupos. El ARI evalúa la identidad de hablante
contra la anotación de AMI y VoxConverse. La recompensa de clips es el $z$ de la fracción de
anotadores que eligió cada segundo en TVSum y SumMe. G-versión verifica la existencia y el
presupuesto de cada lista por video y formato. Las diferencias se reportan con el error
estándar por video (o por reunión), y el umbral de validación es $z > 2$ salvo que el
pre-registro indique otro.


\end{document}
